\documentclass[11pt,a4paper]{article}
\usepackage[T1]{fontenc}
\usepackage[utf8]{inputenc}
\usepackage[margin=25mm]{geometry}
\usepackage{amsmath,amssymb}
\usepackage{mathptmx}
\usepackage{booktabs}
\usepackage{graphicx}
\usepackage{placeins}
\usepackage[round,authoryear]{natbib}
\usepackage{microtype}
\usepackage{hyperref}
\hypersetup{colorlinks=true,linkcolor=black,citecolor=black,urlcolor=blue,
  pdftitle={Termination Constraints and Dynamic Ratio Error in Multistage Differential Resistor-Capacitor Networks}}
\providecommand{\doi}[1]{\href{https://doi.org/#1}{\nolinkurl{doi:#1}}}
\setlength{\emergencystretch}{1.5em}
\clubpenalty=10000
\widowpenalty=10000
\displaywidowpenalty=10000
\raggedbottom
\title{Termination Constraints and Dynamic Ratio Error in Multistage Differential Resistor-Capacitor Networks}
\author{}
\date{}
\begin{document}
\maketitle
\begin{abstract}
We derived termination-dependent cancellation conditions and dynamic ratio-error expressions for grounded differential resistor-capacitor ladders with any finite number of stages. Two recursively computed polynomials describe the open-circuit voltage transfer and output impedance, and a successive extraction test gives necessary and sufficient conditions for positive-component recovery. For a fixed reference branch with $N$ stages, open-circuit matching leaves $N$ local degrees of freedom, one specified positive resistive termination leaves $N-1$, and matching complete transfers at two distinct known terminations requires identical corresponding components. A polynomial-kernel expression extends the dimension calculation to rational loads. For constant resistance, equal direct-current-normalized transfers reject excitation changes exactly in a buffered ratio readout. A changing resistance instead produces a current-weighted history average. A moment expansion separates resistance tracking from the signed current-slew contribution associated with channel-delay mismatch. An 84-case AC matrix spanning one to twelve stages was checked against independent nodal solutions. Eight additional three- and four-stage transient cases, including two time-step refinements, verified both the history dependence and the mismatch term. A load-shifted second-order pair also produced a common-mode step pulse with a calculated unit-step peak of $9.197601$~mV/V and a $20~\mu\mathrm s$ asymptotic time constant. The results connect admissible component families, terminal loading and dynamic fidelity within a common network representation.
\end{abstract}
\noindent\textbf{Keywords:} differential RC filter; common-mode rejection; resistive termination; output impedance; component recovery.


\section{Introduction}

A differential filter processes two input voltages before forming the output difference. Unequal processing converts part of a common input into differential output. \citet{casas1996basics} described matching the transfer functions of noncoupled paths and coupling single-ended filters through a floating common point as two approaches to common-mode rejection. The grounded ladders considered here follow the noncoupled-path configuration.

The relevant matching quantities depend on the measurement frequency and the connected load. \citet{hornero2021cm} showed that component mismatch can produce gain peaking and positive phase shifts in the common-mode response. For a passive ladder, the connected load also draws current through the output impedance. A component combination selected from the open-circuit voltage transfer therefore requires an output-impedance description for predicting the loaded response.

In a buffered resistive-sensor interface, the sample voltage is proportional to $I(t)R(t)$ and the reference voltage to $I(t)R_s$. When resistance changes during the observation window, filtering before forming the ratio produces a weighted history dependence. With mismatched channel responses, changing current can also affect a constant-resistance reading. A matching specification must therefore identify both the terminal conditions for the linear network and the dynamic quantity that the ratio readout estimates.

The objective of this study is to quantify the component freedom available for matching multistage grounded ladders under specified terminations and connect the matching criteria to dynamic ratio measurements. We used classical ladder extraction to express positive realizations in coefficient coordinates, then classified raw and calibrated matching under terminal constraints. Exact constructions and ideal-element circuit simulations checked the load dependence. A convolution analysis separated constant-resistance excitation rejection, resistance-history weighting and the first-order effect of channel-delay mismatch.

\begin{figure}[htbp]
\centering
\includegraphics[width=\linewidth]{figures/Fig_11_en.pdf}
\caption{Signal flow for the buffered ratio model. The highlighted block is the passive network analyzed by the coefficient equations; loads and probe admittances enter at the network output plane. The sensor changes the upstream source amplitude.}
\label{fig:system-plane}
\end{figure}


\section{Related work and the design question}
\label{sec:related-work}

\subsection{Differential transfer functions and impedance balance}
Differential measurement requires both a useful differential response and a small common-to-differential response. Pall\`as-Areny and Webster derived differential-amplifier models with finite differential and common-mode gains and checked the predicted amplitude and phase experimentally \citep{pallasareny1991cmrr}. The companion analysis described cascaded differential stages using four transfer functions and stated the assumption that a stage does not load the preceding stage \citep{pallasareny1991cascaded}. For passive ladders connected directly to a receiver, the receiver impedance enters the branch transfer function. Loading therefore belongs in the matching equations from the outset.

The complex gain carries information needed to compare cancellation conditions. Koide used the complex differential-to-common-mode gain ratio to explain why deliberate imbalance can compensate an active device's common-mode error \citep{koide1997redef}. Gasulla and co-workers showed that common-mode gain and common-to-differential gain are needed to describe transient behavior and stability in fully differential circuits \citep{gasulla2002cmresp}. Hornero and co-workers subsequently analyzed how branch mismatch creates frequency-dependent common-mode responses, including band-pass behavior and phase shifts \citep{hornero2021cm}. The complex-response analyses motivate reporting the separate complex responses over frequency. The magnitude of a rejection ratio alone cannot identify which branch coefficients must agree.

\subsection{Topology, termination, and component matching}
Casas and Pall\`as-Areny analyzed uncoupled and coupled analog differential filters, including higher-order extensions \citep{casas1996basics}. Matching the transfers of uncoupled branches provides cancellation; changing the interconnection can instead produce cancellation through circuit topology. The floating common-point construction establishes an important reference case for topology-based rejection. The circuit considered here consists of two grounded two-stage ladders whose output terminals may carry equal resistive loads and an output-bridging capacitor. The component conditions must consequently be derived for those connections.

Impedance balance also determines the behavior of practical balanced interfaces. Whitlock and Floru showed how source and receiver common-mode impedance imbalances interact and described an input circuit with high common-mode impedance \citep{whitlock2004balancedinput}. The relevant design question for the present ladder is whether two branches with identical open-circuit voltage transfer continue to agree after termination. An open-circuit transfer fixes a voltage ratio, while the output impedance fixes how an attached load changes that ratio. Retaining both functions makes the termination dependence explicit and permits a direct comparison between open-circuit matching and matching at specified finite loads.

\subsection{Network realization and tolerance analysis}
Guillemin addressed RC-network synthesis; Fialkow and Gerst studied RC-ladder transfer functions and general two-terminal-pair RC networks \citep{guillemin1949rcsynthesis,fialkow1951rcladder,fialkow1952rcnetworks}. Wang and co-workers considered passive two-port RC-ladder realization with a specified gain \citep{wang2017rcladder}. The cited synthesis works establish the synthesis context for recovering positive components from transfer coefficients. The present calculation also retains the output-impedance numerator and asks how specified terminations change the set of components that match a fixed reference branch. \emph{Draft note: a proposition-level comparison with the four synthesis papers remains to be completed.}

Component realization and sensitivity must be considered together. Moschytz discussed the trade-offs among sensitivity, component spread, and component tolerance in active-filter design \citep{moschytz0radeffs}. Director and Rohrer developed a generalized adjoint-network framework for obtaining network sensitivities \citep{director1969the}. Endres and Kirkpatrick examined how differential loop-filter mismatch introduces additional transients in fast-settling phase-locked loops \citep{endres0ensitivi}. The sensitivity and transient studies support separating an exact nominal cancellation condition from the effects of component perturbations and from a settling-time specification. Each question requires a specified response function and evaluation criterion.

\subsection{Scope of the present contribution}
The present study derives the complete coefficient conditions for equality of the loaded transfers of two positive, grounded, two-stage RC branches. For a fixed reference branch, it constructs the admissible component family for an open circuit and for one positive resistive termination, and proves component uniqueness when equality holds at two distinct positive resistive terminations. It then examines how a capacitor connected only between the two outputs changes the separate common-to-differential and differential-to-differential gains. Analytical constructions are checked against native AC-simulation outputs.

The cancellation statements concern deterministic linear transfer functions. Serrano-Finetti and co-workers showed that common-mode electronic noise requires a source and cross-correlation model in addition to a deterministic common-mode rejection ratio \citep{serranofinetti2019cmnoise}. A noise-performance claim would therefore require the locations and correlations of the relevant noise sources. The present results supply the nominal network conditions and the termination dependence needed before that further calculation.

\section{Network model and verification method}

\subsection{Signal paths and measurement plane}
\label{sec:model-plane}

Each signal path contained two series resistors, \(R_1\) and \(R_2\). Capacitor \(C_1\) connected the resistor junction to ground, and capacitor \(C_2\) connected the output node to ground. All four component values were positive. Ideal voltage sources drove the two inputs, and the initial conditions were zero. The output was the voltage difference between the two capacitor \(C_2\) nodes. Figure~\ref{fig:circuit} shows the signal paths and termination plane.

We defined the differential and common-mode inputs as

\begin{equation}\label{eq:1}
v_d=v_p-v_n,\qquad v_c=\frac{v_p+v_n}{2}.
\end{equation}

The common-mode-to-differential gain was \(G_{DC}\) = \(v_{d,\mathrm{out}}/v_{c,\mathrm{in}}\) with zero differential input. The differential gain was \(G_{DD}\) = \(v_{d,\mathrm{out}}/v_{d,\mathrm{in}}\) with zero common-mode input. The common-mode rejection ratio was \(\lvert G_{DD}/G_{DC}\rvert\). For independent paths with transfer functions \(H_P\) and \(H_N\),

\begin{equation}\label{eq:2}
G_{DC}=H_P-H_N,\qquad
G_{DD}=\frac{H_P+H_N}{2}.
\end{equation}

Cancellation required \(G_{DC}\) to vanish as a rational function of frequency.

\begin{figure}[htbp]
\centering
\includegraphics[width=0.86\linewidth]{figures/Fig_01_circuit.pdf}
\caption{Two grounded second-order RC branches with output terminations. Each ladder contains two series resistors and two shunt capacitors. The optional capacitor $C_d$ bridges the outputs; the bridge route is dotted blue. Both inputs are driven by ideal voltage sources referenced to ground.}
\label{fig:circuit}
\end{figure}


\subsection{Coefficient representation and load conditions}
\label{sec:coefficients}

We normalized the components by \(R_*\) and \(C_*\) and defined

\begin{equation}\label{eq:3}
r_i=R_i/R_*,\qquad q_i=C_i/C_*,\qquad x=sR_*C_*.
\end{equation}

The four branch coefficients were

\begin{equation}\label{eq:4}
\begin{aligned}
a&=r_1q_1+(r_1+r_2)q_2,\\
b&=r_1r_2q_1q_2,\\
c&=r_1+r_2,\\
d&=r_1r_2q_1.
\end{aligned}
\end{equation}

The open-circuit voltage transfer \(H_0\) and output impedance \(Z_0\) followed from the two node equations:

\begin{equation}\label{eq:5}
H_0(x)=\frac{1}{1+ax+bx^2},\qquad
Z_0(x)=R_*\frac{c+dx}{1+ax+bx^2}.
\end{equation}

A resistance \(R_L\) connected from the output to ground gave the normalized conductance \(g=R_*/R_L\). The loaded transfer was

\begin{equation}\label{eq:6}
H_g(x)=\frac{H_0(x)}{1+Z_0(x)/R_L}
      =\frac{1}{1+gc+(a+gd)x+bx^2}.
\end{equation}

At each load condition, both paths had the same specified resistance \(R_L\). We held the reference branch P fixed when counting the admissible component values of the matching branch N. We compared complete transfer functions at each conductance g. A comparison at two conductances consequently involved two frequency-dependent transfer functions for each path.

\subsection{Exact and numerical verification}

We checked the coefficient identities and example component sets with exact rational arithmetic. The circuit simulations used LTspice 26.1.1, \(R_*=1\,\mathrm{k\Omega}\) and \(C_*=10\,\mathrm{nF}\). The frequency sweep contained 145 points from 10 Hz to 10 MHz, with 24 points per decade. Common-mode excitation applied 1 V to both inputs. Differential excitation applied +0.5 V and -0.5 V to the two inputs.

The 37 simulation cases covered first-order controls, second-order branches, different component-matching strategies, finite source resistance, symmetric and asymmetric resistive terminations, and output bridging capacitance. We independently constructed the nodal equations from the stored netlists and compared the complex node voltages with the simulator outputs. The analytical load-classification results used the ideal-source and equal-resistive-load assumptions in Sections~\ref{sec:model-plane} and~\ref{sec:coefficients}.


\FloatBarrier

\section{Cancellation conditions for an arbitrary number of stages}
\label{sec:general-order}

\subsection{A recursive representation of the ladder}

Consider a branch with $N\geq1$ stages, each containing a positive series resistance $r_j$ followed by a positive shunt capacitance $q_j$. The component normalization and frequency variable are those in Eq.~\eqref{eq:3}. The two-stage circuit in Figure~\ref{fig:circuit} is one instance of this class. With current directed from the network into the terminal load, the chain matrix of stage $j$ is
\begin{equation}\label{eq:general-matrix}
M_j(x)=\begin{bmatrix}1+r_jq_jx&r_j\\q_jx&1\end{bmatrix}.
\end{equation}
Let $(A_j,B_j)$ be the first row of $M_1\cdots M_j$. Appending the last stage gives
\begin{equation}\label{eq:general-forward}
B_j=r_jA_{j-1}+B_{j-1},\qquad
A_j=A_{j-1}+xq_jB_j,
\quad A_0=1,\ B_0=0.
\end{equation}
Induction in Eq.~\eqref{eq:general-forward} gives $\deg A_N=N$, $\deg B_N=N-1$ and $A_N(0)=1$. The leading coefficients are positive. Each stage matrix and the product have determinant one; the first-row polynomials have no nonconstant common factor.

The open-circuit transfer and source-shorted output impedance follow from the terminal equation. A common resistive load on each branch then gives the loaded transfer below, the same termination decomposition used in classical synthesis \citep{guillemin1949rcsynthesis}:
\begin{equation}\label{eq:general-transfer}
H_0=\frac1{A_N},\qquad
\frac{Z_0}{R_*}=\frac{B_N}{A_N},\qquad
H_g=\frac1{A_N+gB_N}.
\end{equation}
For $N=2$, $A_2=1+ax+bx^2$ and $B_2=c+dx$, recovering Eq.~\eqref{eq:5}. Retaining $B_N$ exposes the termination information absent from the open-circuit voltage transfer.

\subsection{Classical ladder extraction and coefficient coordinates}

With the source shorted, $B_N/A_N$ is the dimensionless driving-point impedance seen from the output. The driving-point impedance admits the following ladder continued fraction \citep{guillemin1949rcsynthesis,lee1978chain}
\[
z_j = \left[q_jx+\left(r_j+z_{j-1}\right)^{-1}\right]^{-1},
\qquad z_0=0,\quad z_N=B_N/A_N.
\]
The successive extraction below expresses continued-fraction ladder synthesis \citep{guillemin1949rcsynthesis,lee1978chain} in the $(A_j,B_j)$ coordinates, specialized to single-resistance/single-capacitance arms; recursive transfer calculations for general-arm ladders are classical as well \citep{fialkow1951rcladder}. We use the extraction to count the admissible component families under specified terminal constraints.

The last-stage values can be recovered from the leading coefficients. Writing $\operatorname{lc}(p)$ for the leading coefficient of polynomial $p$, define
\begin{equation}\label{eq:general-inverse}
\begin{aligned}
q_j&=\frac{\operatorname{lc}(A_j)}{\operatorname{lc}(B_j)},&
A_{j-1}&=A_j-xq_jB_j,\\
r_j&=\frac{\operatorname{lc}(B_j)}{\operatorname{lc}(A_{j-1})},&
B_{j-1}&=B_j-r_jA_{j-1}.
\end{aligned}
\end{equation}
For a positive ladder, these operations undo Eq.~\eqref{eq:general-forward} exactly. At each step, $\deg A_{j-1}=j-1$ and, when $j>1$, $\deg B_{j-1}=j-2$. The final pair is $(A_0,B_0)=(1,0)$. Each ratio is positive and uniquely specified.

Conversely, start with real polynomials satisfying $\deg A_N=N$, $\deg B_N=N-1$ and $A_N(0)=1$. If the reductions satisfy the stated highest-power constraints, every recovered component is positive, and the terminal pair is $(1,0)$, reversing the operations reconstructs the original polynomials from a positive ladder. The recovery conditions give a finite necessary and sufficient realization test for this topology. A failed recovery check rejects the candidate coefficient pair.

The forward map from the $2N$ positive components to the $2N$ coefficients of $A_N$ and $B_N$ is polynomial. The inverse map in Eq.~\eqref{eq:general-inverse} is rational with nonzero denominators at every positive realization. Strict positivity and the nonzero leading coefficients persist in a sufficiently small neighborhood. The maps are smooth inverses there, making the coefficient values valid local coordinates for the component family.

\subsection{The number of matching degrees of freedom}

Hold one reference branch fixed. Open-circuit matching fixes the $N$ nonconstant coefficients of $A_N$, leaving $N$ local degrees of freedom in the positive realization domain. Multiple positive realizations of one transfer function are classical; the present count concerns the component family that matches a fixed reference branch within this topology. Uniform impedance scaling changes $B_N$ by one scalar factor while preserving $A_N$. For $N\geq2$, that scaling follows one direction within a larger matching family.

At a specified $g>0$, raw matching, $H_{g,P}=H_{g,N}$, fixes all coefficients of $D_g=A_N+gB_N$. Termination previously entered realizability as a restriction on the attainable gain range \citep{fialkow1951rcladder,fialkow1952rcnetworks}; the count below instead concerns the family that matches a given reference branch. The constant coefficient fixes $B_N(0)$ through $D_g(0)=1+gB_N(0)$. The remaining $N$ coefficients fix the combinations of $A_N$ and $B_N$. The $N+1$ independent linear constraints leave $N-1$ local degrees of freedom, including zero for $N=1$.

Direct-current calibration instead compares $G_g=D_g(0)/D_g$. Fixing the normalized denominator imposes $N$ independent constraints and leaves $N$ local degrees of freedom. In coefficient coordinates the normalized coefficient of $x^N$ can be varied through $A_N$ independently of the lower coefficients. The constraint map therefore has rank $N$. The gain freedom used here has a classical counterpart in specified-gain synthesis, where the poles are fixed and the gain constant varies over a realizable interval \citep{fialkow1951rcladder,paige1960maximum,wang2017rcladder}. In the present matching problem that freedom is an ambiguity of the partner branch after calibration, with consequences for the uncalibrated differential output. For $N=1$, the positive family
\[
r>0,\qquad q=\frac{r+10}{11r}
\]
has the same normalized response at $g=0.1$ as $(r,q)=(1,1)$, although the raw direct-current gain varies with $r$.

For $N=2$, branch P = $(1,1,1,1)$ and the positive branch
\[
\left(2,2,\frac{217+\sqrt{6769}}{288},\frac{217-\sqrt{6769}}{480}\right)
\]
have identical normalized responses at $g=0.1$. The two raw denominators differ by the constant factor $7/6$, giving a raw CMRR of $20\log_{10}(13/2)=16.26$~dB at every frequency. Calibration therefore permits an additional matching direction that does not cancel the raw differential output.

\begin{table}[htbp]
\centering
\caption{Local matching dimensions under resistive terminations, for a fixed positive $N$-stage reference branch. Raw matching includes absolute gain; normalized matching compares $H_g/H_g(0)$.}
\label{tab:matching-conventions}
\begin{tabular}{lll}
\toprule
Termination information & Raw & Normalized \\
\midrule
Open circuit & $N$ & $N$ \\
One specified $g>0$ & $N-1$ & $N$ \\
Two distinct specified $g_1,g_2>0$ & $0$ & $0$ \\
\bottomrule
\end{tabular}
\end{table}

At two distinct positive conductances, put $U_i=1/H_{g_i}$. The two complete transfer functions determine
\begin{equation}\label{eq:general-two-load}
B_N=\frac{U_2-U_1}{g_2-g_1},\qquad
A_N=\frac{g_2U_1-g_1U_2}{g_2-g_1}.
\end{equation}
Two branches that agree at both loads consequently have identical $A_N$ and $B_N$. The stepwise uniqueness of the classical ladder extraction then requires identical corresponding components \citep{guillemin1949rcsynthesis}. The two-load reconstruction applies to every finite $N$ in the stated ladder class. It uses complete complex transfer functions, including absolute gain, at two loads.

Two normalized resistive-load responses also identify the components exactly. Write $D_i(g_k)=c_kD_{\rm ref}(g_k)$ for a candidate branch. The coefficient of $x^N$ is independent of $g_k$ and is nonzero; hence $c_1=c_2=c$. Subtracting the two load equations gives $B_i=cB_{\rm ref}$ and then $A_i=cA_{\rm ref}$. Since $A_i(0)=A_{\rm ref}(0)=1$, $c=1$. The normalized two-load proof uses exact complete responses. A stable noisy-data estimator additionally requires an uncertainty model.

\subsection{The load type also determines the matching family}

A known rational load can be written as a dimensionless admittance $g(x)=P(x)/Q(x)$, where $P$ and $Q$ are nonzero coprime real polynomials. The loaded response is defined provided $QA_N+PB_N$ is not identically zero. Equality of the two branch transfers gives
\begin{equation}\label{eq:rational-kernel}
H_g=\frac{Q}{QA_N+PB_N},\qquad
Q\Delta A+P\Delta B=0.
\end{equation}
Here $\Delta$ denotes branch P minus branch N. Coprimality implies $\Delta A=PS$ and $\Delta B=-QS$ for a polynomial $S$. The highest-power bounds and $\Delta A(0)=0$ restrict $S$ to
\begin{equation}\label{eq:rational-dimension}
\begin{aligned}
m&=\min\{N-\deg P,\ N-1-\deg Q\},\\
d_g&=\max\{0,m+1-\epsilon\},\qquad
\epsilon=\begin{cases}1,&P(0)\ne0,\\0,&P(0)=0.\end{cases}
\end{aligned}
\end{equation}
For $m<0$, the only allowable polynomial is $S=0$. Otherwise, $S$ has $m+1$ coefficients, and a nonzero $P(0)$ fixes the constant coefficient of $S$ to zero. Intersecting this linear family with the open positive realization domain preserves the local dimension $d_g$.

Table~\ref{tab:load-dimensions} applies Eq.~\eqref{eq:rational-dimension} to common loads. The open circuit, for which $P$ is identically zero, is listed separately. A pure capacitor changes the terminal capacitance and retains $N$ degrees of freedom, whereas a resistor also constrains the zero-frequency gain and reduces the count to $N-1$.
\begin{table}[htbp]
\centering
\caption{Local raw matching dimension for a fixed reference branch under one specified load, including absolute gain. All displayed load parameters are positive. Direct-current-normalized matching is treated separately in Table~\ref{tab:matching-conventions}.}
\label{tab:load-dimensions}
\begin{tabular}{lll}
\toprule
Load & Dimensionless admittance & Dimension \\
\midrule
Open circuit & $0$ & $N$ \\
Resistor & $g_0$ & $N-1$ \\
Capacitor & $\alpha x$ & $N$ \\
Inductor & $\beta/x$ & $\max(N-2,0)$ \\
Series RC & $\alpha x/(1+\gamma x)$ & $N-1$ \\
Parallel RC & $g_0+\alpha x$ & $N-1$ \\
\bottomrule
\end{tabular}
\end{table}

Two distinct known rational admittances also determine the component tuple under raw matching. For $g_k=P_k/Q_k$, multiply $Q_k\Delta A+P_k\Delta B=0$ by the other denominator and subtract. The result is $(P_1Q_2-P_2Q_1)\Delta B=0$. Nonidentical load functions mean $P_1Q_2-P_2Q_1\not\equiv0$, which gives $\Delta B=0$ and then $\Delta A=0$. Each known nonzero $Q_k$ permits cancellation as a rational function, including at isolated zeros by continuation. The loads in Table~\ref{tab:load-dimensions} are time-invariant rational admittances; a periodically switched input requires the separate sampling model introduced in Section~\ref{sec:dynamic}.

\subsection{Load variation and recovery conditioning}

For independent branches, let $D_i(g)=A_i+gB_i$. The exact common-mode-to-differential gain is
\begin{equation}\label{eq:general-leakage}
G_{DC}=-\frac{\Delta A+g\Delta B}{D_P(g)D_N(g)}.
\end{equation}
If the branches match at $g_0>0$, then $\Delta A=-g_0\Delta B$ and $\Delta B(0)=0$. Under a different common resistive load,
\begin{equation}\label{eq:general-load-shift}
G_{DC}(g)=-\frac{(g-g_0)\Delta B}{D_P(g)D_N(g)}.
\end{equation}
The zero-frequency leakage remains zero, while the remaining polynomial $\Delta B$ controls the frequency-dependent mismatch. The second-order load-shift formula is the case $\Delta B=x\Delta d$.

Exact identifiability and measurement conditioning require separate checks. Equation~\eqref{eq:general-two-load} divides by the load separation, and a small response perturbation gives $\delta(1/H)\simeq-\delta H/H^2$. Closely spaced loads and strongly attenuated responses amplify measurement errors. Experimental recovery therefore needs a stated frequency band and uncertainty model in addition to the algebraic uniqueness result.
For the dimensionless coefficient map $(r_1,\ldots,r_N,q_1,\ldots,q_N)\mapsto(A_{N,1},\ldots,A_{N,N},B_{N,0},\ldots,B_{N,N-1})$, the unscaled Jacobian at unit components has 2-norm condition numbers $15.6$, $1.08\times10^3$, $9.30\times10^4$, $8.67\times10^6$, $8.40\times10^8$ and $8.35\times10^{10}$ for $N=2,4,6,8,10,12$, respectively. The condition numbers concern the stated coefficient normalization and change under coordinate rescaling. They quantify local amplification in coefficient recovery; they are not a bound on the performance of a particular response-fitting estimator.
\FloatBarrier

\section{Cancellation conditions and component recovery}

\subsection{Positive-component recovery}

The four coefficients determined the components through

\begin{equation}\label{eq:7}
q_2=\frac{b}{d},\qquad
r_2=\frac{d^2}{ad-bc},\qquad
r_1=c-r_2,\qquad
q_1=\frac{d}{r_1r_2}.
\end{equation}

Figure~\ref{fig:positive-region} illustrates the positive-component domain at fixed $a$ and $b$. A coefficient tuple admitted positive components exactly when

\begin{equation}\label{eq:8}
b>0,\qquad d>0,\qquad ad-bc>0,\qquad
c(ad-bc)>d^2.
\end{equation}

The first two inequalities made \(q_2\) positive. The third made \(r_2\) positive, and the fourth made \(r_1\) positive. The recovered \(q_1\) was then positive. Substitution of Eq.~\eqref{eq:7} into Eq.~\eqref{eq:4} recovered the original coefficients. Conversely, positive components satisfied Eq.~\eqref{eq:8}, including the identity \(ad-bc=r_1^2r_2q_1^2\). The mapping was therefore one-to-one over the positive-component domain.

\begin{figure}[htbp]
\centering
\includegraphics[width=0.67\linewidth]{figures/Fig_02_positive_region.pdf}
\caption{Positive-component region for fixed open-circuit coefficients $a=3$, $b=1$. The shaded interior satisfies $d>0$ and $c(3d-c)>d^2$; boundary lines are excluded. P and Q are distinct realizations of the same open-circuit transfer. Coordinates are normalized coefficients.}
\label{fig:positive-region}
\end{figure}


\subsection{Open-circuit cancellation}

For open outputs, \(G_{DC}\) in Eq.~\eqref{eq:2} vanished exactly when

\begin{equation}\label{eq:9}
a_P=a_N,\qquad b_P=b_N.
\end{equation}

Fixing a and b left c and d free within Eq.~\eqref{eq:8}. Equal open-circuit transfer therefore admitted a two-parameter component family. Uniform impedance scaling, \(r_i\rightarrow kr_i\) and \(q_i\rightarrow q_i/k\), preserved a and b and supplied one construction within that family.

A second construction demonstrated the additional freedom. In the component order \((r_1,r_2,q_1,q_2)\), branches P = (1,1,1,1) and Q = (1,4,1/2,1/2) both gave a = 3 and b = 1. Different Q/P resistor ratios of 1 and 4 ruled out a uniform impedance-scaling factor. The output-impedance coefficients were \((c,d)\) = (2,1) for P and (5,2) for Q.

\subsection{Cancellation at one positive load conductance}

For a specified \(g>0\), equality of Eq.~\eqref{eq:6} required

\begin{equation}\label{eq:10}
c_P=c_N,\qquad b_P=b_N,\qquad
 a_P+gd_P=a_N+gd_N.
\end{equation}

The constant numerators made equality of the transfer functions equivalent to equality of the denominator coefficients. The three conditions in Eq.~\eqref{eq:10} left one free coefficient in the positive-component domain. Choosing d fixed a through the prescribed value of a + gd, and Eq.~\eqref{eq:7} then recovered the components.

Branches P and T illustrated this family. Branch T = (18/19,20/19,361/180,1/2) gave \((a,b,c,d)\) = (2.9,1,2,2). At g = 0.1, P and T both had the denominator \(1.2+3.1x+x^2\). At g = 0.2, the denominators became \(1.4+3.2x+x^2\) for P and \(1.4+3.3x+x^2\) for T. A load change separated the responses that had matched at the first termination.

\subsection{Cancellation at two distinct positive load conductances}

Suppose the two branches had equal complete transfer functions at \(g_1\) and \(g_2\), with \(g_1>0\), \(g_2>0\) and \(g_1\ne g_2\). Equality of the constant and quadratic denominator coefficients gave \(c_P\) = \(c_N\) and \(b_P\) = \(b_N\). Equality of the linear coefficients gave

\begin{equation}\label{eq:11}
\Delta a+g_1\Delta d=0,\qquad
\Delta a+g_2\Delta d=0,
\end{equation}

where \(\Delta\) denoted the difference between branch coefficients. Subtraction gave \(\Delta d=0\), followed by \(\Delta a=0\). All four coefficients were equal. The inverse mapping in Eq.~\eqref{eq:7} consequently required equal corresponding component values.

For the fixed reference branch, the two-load requirement reduced the admissible matching family from one free coefficient to a single component tuple. Identical components also satisfied the requirement directly, establishing necessity and sufficiency for the stated ladder class.


\subsection{Leakage caused by load changes and component deviations}

Let $D_i(g)=1+gc_i+(a_i+gd_i)x+b_ix^2$ denote the denominator of branch $i$. The exact difference between two branch transfers is
\begin{equation}\label{eq:load-leakage}
G_{DC}=\frac{D_N(g)-D_P(g)}{D_P(g)D_N(g)}.
\end{equation}
For branches matched at a positive conductance $g_0>0$, define $\Delta d=d_P-d_N$. Equation~\eqref{eq:10} then gives
\begin{equation}\label{eq:load-shift}
G_{DC}(g)=-\frac{(g-g_0)\Delta d\,x}{D_P(g)D_N(g)}.
\end{equation}
The numerator identifies the joint roles of load change and the remaining output-impedance coefficient. A pair with unequal $d$ retains zero leakage at DC after a common load change, but develops frequency-dependent leakage away from DC. Matching the DC outputs alone consequently leaves the one-load ambiguity unresolved.

For small coefficient differences around a matched branch with transfer $H_g$, the first-order leakage is
\begin{equation}\label{eq:sensitivity}
G_{DC}\simeq-H_g^2\left[g\Delta c+x(\Delta a+g\Delta d)+x^2\Delta b\right].
\end{equation}
Here each difference is taken as branch P minus branch N; the neglected terms are second order in the coefficient deviations at frequencies where the branch denominator is nonzero. The three frequency-dependent weights separate the contributions of the constant, linear and quadratic denominator coefficients. A tolerance allocation can therefore target the loaded coefficient combinations directly.

Unequal external loads introduce an additional conversion mechanism even when the component values of the branches are identical. With conductances $g_P$ and $g_N$,
\begin{equation}\label{eq:load-asymmetry}
G_{DC}=-\frac{(g_P-g_N)(c+dx)}{D(g_P)D(g_N)}.
\end{equation}
Equation~\eqref{eq:load-asymmetry} makes load balance part of the cancellation specification. Equal components give cancellation under a common termination; the two attached load admittances must also agree.

\FloatBarrier

\section{Dynamic ratio error and transient common-mode leakage}
\label{sec:dynamic}

The cancellation conditions above are exact identities between transfer functions, evaluated at every frequency. This section derives what they imply when the measured resistance and the excitation current vary during the observation window. The model remains the buffered front end of Section~\ref{sec:model-plane}: the sample resistance modulates the source amplitude, and the ladder itself is linear and time-invariant.

\subsection{Matched channels: an exact weighted-average identity}
\label{sec:dynamic-matched}

Let the two loaded branch transfers be equal after direct-current gain normalization, $G_P=G_N=G$ with $G(0)=1$ and impulse response $g(t)$, and form the ratio readout
\begin{equation}\label{eq:ratio-readout}
\hat R(t)=R_s\,\frac{(g\!\ast\! IR)(t)}{(g\!\ast\! I)(t)} ,
\end{equation}
where $I(t)$ is the excitation current and $R(t)$ the sample resistance. Because $R(t)$ is constant with respect to the convolution variable,
\begin{equation}\label{eq:weighted-average}
\hat R(t)-R(t)=\frac{\bigl(g\!\ast\! I\,[R(\cdot)-R(t)]\bigr)(t)}{(g\!\ast\! I)(t)} .
\end{equation}
Equation~\eqref{eq:weighted-average} is exact for every waveform of $I(t)$, including steps and square-wave ripple with fast edges: \emph{any} excitation-current fluctuation cancels from the ratio when the loaded transfers match. The remaining error is a weighted average of the past resistance values. If $|{\rm d}R/{\rm d}t|\le L_R$ and $|I|\le I_{\max}$, then
\begin{equation}\label{eq:lag-bound}
|\hat R(t)-R(t)|\;\le\;\frac{L_R\,I_{\max}\,M_1}{\bigl|(g\!\ast\! I)(t)\bigr|},
\qquad
M_1=\int_0^\infty \!\!\tau\,|g(\tau)|\,{\rm d}\tau .
\end{equation}
The bound uses the first absolute moment because the impulse response of an $N\ge2$ ladder is not nonnegative; residues alternate in sign although all poles are real and negative. For the ladder class, $M_1\le\sum_k |r_k|/|p_k|^2$ in terms of the poles and residues of $1/D_g$, which the coefficients of Section~\ref{sec:general-order} supply directly.

\subsection{Mismatched channels: the current-slew cross term}
\label{sec:dynamic-mismatch}

When the loaded transfers differ, $G_P\ne G_N$, the direct-current normalization removes the static ratio error, and the ratio error acquires a term driven by the current slew rate. Writing $\tau=D'(0)/D(0)$ for the direct-current group delay of the loaded transfer (in $x$ units, $\tau=(A_1+gB_1)/(1+gB_0)$, scaled by $R_*C_*$), the first-order expression is
\begin{equation}\label{eq:mismatch-dynamic}
\hat R-R \;=\; \text{lag term of Eq.~\eqref{eq:lag-bound}} \;+\; R\,\frac{\dot I}{I}\,\Delta\tau
\;+\;O\!\left(\ddot I/I,\ \Delta M_1\right),
\qquad \Delta\tau=\tau_P-\tau_N .
\end{equation}
The cross term $R\,(\dot I/I)\,\Delta\tau$ is absent from every static analysis: a pair that cancels common-mode perfectly at each frequency can still convert excitation-current dynamics into a ratio error through the group-delay mismatch of the two channels. For the second-order unit branch at $g=0.1$, a $+10\%$ deviation of the output capacitor gives $\Delta\tau\approx1.67~\mu{\rm s}$; at a relative current slew $\dot I/I=800~{\rm s}^{-1}$ and $R=100~{\rm k}\Omega$ this is $133~\Omega$ of dynamic ratio error.

\subsection{Equivalence of static cancellation and dynamic immunity}
\label{sec:welding}

The static conditions of Section~\ref{sec:general-order} and the dynamic property of Section~\ref{sec:dynamic-matched} are one statement. At the operating load $g$ the following are equivalent:
\begin{enumerate}
\item $D_P=D_N$ (equal coefficient vectors, the static cancellation condition);
\item $g_P=g_N$ (equal impulse responses);
\item the ratio readout~\eqref{eq:ratio-readout} rejects every excitation-current waveform, in the sense of Eq.~\eqref{eq:weighted-average}.
\end{enumerate}
Statements 1 and 2 are equal polynomials of equal degree. Statement 2 implies 3 by the derivation of Eq.~\eqref{eq:weighted-average}; conversely, driving both channels with an impulse-pair input isolates $g_P-g_N$, so 3 implies 2. The matching-family dimensions of Table~\ref{tab:load-dimensions} therefore count not only static common-mode cancellation but also the families inside which the dynamic ratio readout is exactly immune to excitation dynamics.

\subsection{Transient common-mode leakage after a load change}
\label{sec:cm-pulse}

For the load-shift family matched at $g_0$ and observed at $g$, Eq.~\eqref{eq:general-load-shift} gives $G_{DC}=-(g-g_0)\Delta B/(D_P D_N)$ with $\Delta B(0)=0$. The zero-frequency leakage vanishes, yet the response is band-limited, so a common-mode \emph{step} of amplitude $V_c$ produces a finite differential pulse
\begin{equation}\label{eq:cm-pulse}
y(t)=V_c\,\mathcal L^{-1}\!\left[\frac{G_{DC}(s)}{s}\right],
\qquad
\|y\|_2^2=V_c^2\sum_{k,j}\frac{r_k r_j}{-(p_k+p_j)} ,
\end{equation}
with the residue formula for the pulse energy and a peak bounded by $\|g_{DC}\|_1$; the slowest pole of $D_P D_N$ governs the recovery time. For the P/T branches of Table~\ref{tab:branches} matched at $g_0=0.1$ and observed at $g=0.2$, the closed-form step response peaks at $9.2\times10^{-3}\,V_c$ after $28~\mu{\rm s}$ and decays with a $19~\mu{\rm s}$ tail (all four poles real, between $-2.69$ and $-0.52$ in $x$ units).

\subsection{Transient verification}
\label{sec:dynamic-verif}

Four transient simulations checked these predictions with \(R_*=1~\mathrm{k\Omega}\), \(C_*=10~\mathrm{nF}\), and a \(10~\mathrm{k\Omega}\) operating load. With matched branches, a \(5~\mathrm{kHz}\) square-wave current ripple of \(\pm20\%\) with \(10~\mathrm{ns}\) edges and a resistance ramp from \(10~\mathrm{k}\Omega\) to \(120~\mathrm{k}\Omega\) in \(10~\mathrm{ms}\) gave a maximum ratio error of \(320.93~\Omega\), inside the bound of Eq.~\eqref{eq:lag-bound} with a \(1.08\times\) margin; the transient solver and an independent fourth-order reference solution agreed to six digits. With an output capacitor deviating by \(+10\%\) and a \(1~\mathrm{kHz}\) triangular current, the measured mean error of \(121.9~\Omega\) matched the cross-term prediction \(R(\dot I/I)\Delta\tau=133.3~\Omega\) within nine percent, and the joint scenario amplified the error to \(4.9\times\) the matched baseline. The common-mode step pulse of Section~\ref{sec:cm-pulse} agreed with the closed form to a relative difference of \(1.2\times10^{-7}\).

\section{Simulation results and output coupling}

\subsection{General-order verification and model consistency}

The 84-case matrix covered open and resistively loaded branches, several pairing rules, nonuniform same-transfer realizations, fixed-load matching families, reactive loads, source impedance, output coupling and two-load recovery. Orders ranged from one to twelve. With the documented 1~m$\Omega$ inductor series resistance included in the independent model, the largest absolute complex node-voltage discrepancy was $1.68\times10^{-14}$~V.

Against the originally intended ideal-inductor reference, two cases exceeded the unchanged $10^{-7}$~V gate, with errors of $4.54\times10^{-7}$ and $3.57\times10^{-7}$~V. Fresh runs specifying $R_{\rm ser}=0$ agreed with the ideal reference to less than $1.8\times10^{-15}$~V; explicit 1~m$\Omega$ runs reproduced the default-model behavior. The discrepancy arose from different element models. Both original failures and the diagnostic records were retained.

At $10$~kHz, the three-stage same-$A$/different-$B$ construction agreed in the open circuit to a relative difference of $1.4\times10^{-14}$ and separated under a 10~k$\Omega$ load by $3.4\times10^{-4}$. A fixed-load family agreed at $g=0.1$ to $6.7\times10^{-14}$ and differed at $g=0.2$ by $3.2\times10^{-5}$. A capacitive-load matching family agreed under the specified capacitance to $5.3\times10^{-14}$ and differed under resistive loading by $2.2\times10^{-3}$. The comparisons tested complete complex outputs; the following second-order examples make the same mechanisms visible with shorter component tuples.

\subsection{Termination-dependent equivalence}

Table~\ref{tab:branches} collects the three branch constructions. The equalities and inequalities were established algebraically before the frequency sweeps.

\begin{table}[htbp]
\centering
\caption{Normalized component values and branch coefficients.}
\label{tab:branches}
\begin{tabular}{lrrrrrrrr}
\toprule
Branch & $r_1$ & $r_2$ & $q_1$ & $q_2$ & $a$ & $b$ & $c$ & $d$ \\
\midrule
P & 1 & 1 & 1 & 1 & 3 & 1 & 2 & 1 \\
Q & 1 & 4 & $1/2$ & $1/2$ & 3 & 1 & 5 & 2 \\
T & $18/19$ & $20/19$ & $361/180$ & $1/2$ & 2.9 & 1 & 2 & 2 \\
\bottomrule
\end{tabular}
\end{table}

Figure~\ref{fig:termination-transfer} compares the open and loaded P/Q responses. Figure~\ref{fig:output-impedance} shows the output impedances that distinguish the open-circuit-equivalent branches. Figure~\ref{fig:load-specific} compares T with P after each load is applied.



The simulations reproduced open-circuit cancellation for P/Q and cancellation at \(R_L=10\,\mathrm{k\Omega}\) for P/T. Changing \(R_L\) to 5 k\(\Omega\) produced unequal P/T responses. Across the 37 cases, the maximum absolute complex node-voltage difference between LTspice and the independently assembled nodal solution was \(1.14\times10^{-15}\) V.

\begin{figure}[htbp]
\centering
\includegraphics[width=0.67\linewidth]{figures/Fig_03_termination_transfer.pdf}
\caption{Magnitude responses of P and Q under open-circuit and 10~k$\Omega$ shunt-load conditions. Lines are analytical transfer functions; open circles are archived LTspice samples. Open-circuit responses coincide. The normalization is $R_*=1$~k$\Omega$, $C_*=10$~nF.}
\label{fig:termination-transfer}
\end{figure}

\begin{figure}[htbp]
\centering
\includegraphics[width=0.67\linewidth]{figures/Fig_04_output_impedance.pdf}
\caption{Open-circuit output-impedance magnitudes of P, Q and T with the ideal input voltage sources set to zero. Curves are calculated from the four-coefficient representation. P and Q share an open-circuit voltage transfer but have different output impedances.}
\label{fig:output-impedance}
\end{figure}

\begin{figure}[htbp]
\centering
\includegraphics[width=0.67\linewidth]{figures/Fig_05_load_specific_match.pdf}
\caption{Relative magnitude response of T to P at two common shunt loads. The ordinate is $20\log_{10}|H_T/H_P|$; zero indicates equal magnitudes. Lines are analytical responses; open circles are archived LTspice samples. Complete-function equality at 10~k$\Omega$ is established algebraically in the text.}
\label{fig:load-specific}
\end{figure}


\subsection{A bridging capacitor acts on both differential-output gains}

We connected a capacitor \(C_d\) between the two output nodes while retaining the ideal sources and the specified shunt loads. Let \(Z'_P\text{ and }Z'_N\) denote the branch output impedances in parallel with the respective shunt loads, and let \(H_P\) and \(H_N\) denote the loaded transfers before \(C_d\) was added. Solving the two output-node equations gave

\begin{equation}\label{eq:12}
\begin{aligned}
G_{DC}&=\frac{H_P-H_N}{D_d},\\
G_{DD}&=\frac{(H_P+H_N)/2}{D_d},\\
D_d&=1+sC_d(Z'_P+Z'_N).
\end{aligned}
\end{equation}

The shared factor \(D_d\) cancelled from \(G_{DD}/G_{DC}\) wherever the ratio was defined. The bridging capacitor changed both differential-output gains by the same factor and preserved \(G_{DD}/G_{DC}\). Equal loaded branch transfers also preserved \(G_{DC}\) = 0 after the capacitor was connected.

For P/Q with \(R_L=10\,\mathrm{k\Omega}\) and \(C_d=820\,\mathrm{pF}\), the maximum complex-ratio difference was \(5.14\times10^{-12}\) relative to the ratio without the bridging capacitor, evaluated over 145 frequency points. The ratio followed the cancellation of the shared factor in Eq.~\eqref{eq:12}.


Figure~\ref{fig:bridge-gains} displays the changes in the separate gains, while Figure~\ref{fig:bridge-cmrr} compares the corresponding rejection ratios.

\begin{figure}[htbp]
\centering
\includegraphics[width=0.67\linewidth]{figures/Fig_06_bridge_gains.pdf}
\caption{Analytical differential-output gains of the loaded P/Q pair with and without a bridging capacitor. DM denotes unit differential input; CM$\rightarrow$DM denotes unit common-mode input and differential output. Open markers identify analytical samples on the 820~pF curves. Both shunt loads are 10~k$\Omega$.}
\label{fig:bridge-gains}
\end{figure}


\begin{figure}[htbp]
\centering
\includegraphics[width=0.67\linewidth]{figures/Fig_07_bridge_cmrr.pdf}
\caption{Common-mode rejection ratio of the loaded P/Q pair. The analytical curves for $C_d=0$ and 820~pF coincide. Open circles are calculated from archived common-mode and differential-mode LTspice output differences with $C_d=820$~pF. Both shunt loads are 10~k$\Omega$; CMRR is $20\log_{10}|G_{DD}/G_{DC}|$. Absolute gains accompany the ratio in Figure~\ref{fig:bridge-gains}; the independent-branch termination sweep is shown in Figure~\ref{fig:terminal-sweep}.}
\label{fig:bridge-cmrr}
\end{figure}


\FloatBarrier

\section{Discussion}

The operating termination determines the component freedom available for cancellation. A fixed $N$-stage reference branch admits an $N$-dimensional open-circuit matching family and an $(N-1)$-dimensional family at one positive resistive load. Two distinct known terminations fix the complete polynomial pair and therefore the components. The rational-load result further separates the effect of circuit order from the effect of load type.

The coefficient representation gives a constructive design sequence. Specify the source and output measurement planes and the terminal admittance. Select the required transfer coefficients, choose the remaining free coefficients, and apply the recursive positive-realization test in Eq.~\eqref{eq:general-inverse}. Evaluate component and load perturbations on the recovered circuit. The P/Q and P/T constructions illustrate why an open-circuit equivalence or a single-load match must be tested at the intended operating termination.

\citet{casas1996basics} obtained cancellation through floating coupling in a different topology, including higher-order extensions. The grounded ladders here retain shunt current paths, and the pair $(A_N,B_N)$ exposes terminal loading. \citet{hornero2021cm} analyzed frequency-dependent mismatch through the common-mode response. The present representation connects that response to a positive component family and then to the calibrated ratio observable.

Strict component positivity preserves the stated order and the uniqueness result. For example, setting $q_1=0$ in the second-order ladder merges two series resistors. $(1,2,0,3)$ and $(3/2,3/2,0,3)$ yield the same coefficients $(9,0,3,0)$. The inverse then loses the separation of the two resistances. Unknown source resistance similarly combines with the first series resistor unless the source contribution is measured independently.

\subsection{Dynamic fidelity and calibration}

Direct-current calibration changes which transfer functions need to match. A constant-resistance ratio rejects all admissible current histories when the normalized channel responses are equal. Raw common-mode cancellation additionally requires equal absolute gains. The matching dimensions in Table~\ref{tab:load-dimensions} concern the raw response and retain that gain requirement.

For changing resistance, matched channels preserve a weighted history average. Reducing the direct-current group delay $\tau=D'(0)/D(0)$ at the operating load reduces the first-moment tracking bound, while reducing the channel-delay mismatch $\Delta\tau=\tau_P-\tau_N$ (Section~\ref{sec:dynamic-mismatch}) suppresses the additional first-order current-slew term. Tracking delay and channel matching are different design objectives. The three- and four-stage transient examples quantify both effects. Identical channels retained current-dependent history weighting, and capacitor mismatch generated the predicted signed error during excitation changes.

The two-load result supplies an identification strategy for $2N$ component parameters. In an experiment, complex responses under known loads would determine the polynomial pair with uncertainty intervals. Closely spaced loads and strongly attenuated responses can amplify errors in that recovery, as described in Section~\ref{sec:general-order}. Measurement-band selection and conditioning therefore accompany the exact identifiability result.

\subsection{Scope of the evidence and experimental limitations}

The evidence in this paper consists of exact algebra and independent circuit simulation; no physical measurements are included. A bench validation of the termination-dependent families would require measured component values, characterized complex source and probe impedances, separate branch responses under known terminations, common-mode and differential excitation of the pair, repeatability records, and an uncertainty budget sufficient to separate nominal from measured-component predictions.

Dynamic validation additionally requires the two source signals and both filtered outputs on a common time base, with constant-resistance and changing-resistance criteria evaluated separately and the measured waveforms driving the reference model, so that excitation shape and channel mismatch remain distinguishable. Leakage measurements near the instrument floor are bounded by the corresponding detection limit rather than resolved.


\section{Conclusions}

We established termination-dependent cancellation conditions for grounded differential resistor-capacitor ladders, first in full detail at second order and then for an arbitrary number of stages through a recursive chain-matrix representation. For a fixed reference branch, open-circuit cancellation admitted a two-parameter matching family; one positive resistive termination admitted a one-parameter family; and equality of the complete transfer functions at two distinct positive resistive terminations required identical corresponding components. Positive-component recovery conditions and explicit branch constructions made each result directly checkable. The analytical predictions were reproduced in 37 second-order and 84 general-order AC simulation cases, and an output bridging capacitor preserved the common-mode rejection ratio through a shared loading factor in the two differential-output gains. For time-varying measurements, equal loaded transfers were shown to be exactly the condition under which the ratio readout rejects arbitrary excitation-current dynamics; a first-order cross term and a transient common-mode leakage pulse quantified the dynamic cost of termination mismatch.



\FloatBarrier
\bibliographystyle{plainnat}
\bibliography{references}
\end{document}
